\documentclass[french]{article}
\usepackage[utf8]{inputenc}
\usepackage[T1]{fontenc}
\usepackage{lmodern}
\usepackage{geometry}
\usepackage{graphicx}
\usepackage{babel}
\usepackage{amsmath}
\usepackage{amsthm}
\usepackage{amssymb}
\usepackage{hyperref}
\usepackage{stmaryrd}
\usepackage{enumitem}
\usepackage{tcolorbox}
\usepackage{dsfont}
\usepackage{multirow}
\usepackage{etoc}
\usepackage{titlesec}
\usepackage{esint}
\usepackage{cancel}
\usepackage{mathdots}

\setlist[itemize]{label=$\bullet$}


\renewcommand{\P}{\mathbb{P}}
\newcommand{\N}{\mathbb{N}}
\newcommand{\Z}{\mathbb{Z}}
\newcommand{\Q}{\mathbb{Q}}
\newcommand{\R}{\mathbb{R}}
\newcommand{\C}{\mathbb{C}}
\newcommand{\K}{\mathbb{K}}
\renewcommand{\L}{\mathbb{L}}
\newcommand{\U}{\mathbb{U}}
\newcommand{\st}{^*}
\newcommand{\tq}{\middle|}
\newcommand{\inv}{^{-1}}
\newcommand{\E}{\mathbb{E}}
\newcommand{\F}{\mathbb{F}}
\newcommand{\V}{\mathbb{V}}
\renewcommand{\ker}{\text{Ker}}
\newcommand{\im}{\text{Im}}
\newcommand{\card}{\text{Card}}
\newcommand{\Tr}{\text{Tr}}

\title{TD Euclidiens\\ Exercice 39}
\date{}
\begin{document}
\maketitle

\section{Énoncé $\star$}
\begin{enumerate}
\item Montrer que $AB$ est diagonalisable. On pourra utiliser une racine carrée.
\item On suppose $B$ positive et on note respectivement $\lambda_A$ et $\lambda_B$ les plus grands valeurs propres de $A$ et $B$. Montrer que toute valeur propre de $AB$ est inférieure ou égale à $\lambda_A\lambda_B$.
\item Montrer que le résultat de la question précédente reste vrai si l'on suppose seulement $A$ et $B$ positives. 
\end{enumerate}
\section{Solution}
\begin{enumerate}
	\item En effet, avec $C$ la racine carrée de $A$ (donc inversible), on écrit : $$AB=C(CBC)C\inv=CP^T DPC\inv$$ avec le théorème spectral appliqué à $B$.
	\item Soit $X\in\R^n$. On a : $$(X\mid CBCX)=(CX\mid BCX)\leqslant\lambda_B\lVert CX\rVert^2=\lambda_B(CX\mid CX)=\lambda_B(X\mid AX)\leqslant \lambda_A\lambda_B\lVert X\rVert^2$$ Pour un vecteur propre $X$ de $AB$ associé à la valeur propres $\lambda$, on obtient : $$\lambda\lVert X\rVert^2\leqslant\lambda_1\lambda_B\lVert X\rVert^2$$ donc $\lambda\leqslant\lambda_A\lambda_B$.
	\item On utilise le fait que la plus grande valeur propre vaut $\displaystyle\underset{\lVert X\rVert=1}{\max}(X\mid AX)$ puis on raisonne par densité et compacité de la sphère unité.
\end{enumerate}
\end{document}
